\documentclass[10pt,twocolumn,letterpaper]{article}
\usepackage[applications,review]{wacv}

\usepackage{microtype}
\usepackage{graphicx}
\usepackage{amsmath,amssymb}
\usepackage{booktabs}
\usepackage{multirow}
\usepackage{makecell}
\usepackage{caption}
\usepackage{siunitx}
\sisetup{detect-weight=true, detect-family=true}
% Global table compaction.
\renewcommand{\arraystretch}{0.92}
\setlength{\tabcolsep}{4pt}
\setlength{\aboverulesep}{0pt}
\setlength{\belowrulesep}{0pt}
\usepackage{xcolor}
\usepackage{colortbl}
\usepackage[pagebackref,breaklinks,colorlinks]{hyperref}
\usepackage[capitalize,noabbrev]{cleveref}

% Earthy greenish-brown palette (FTP theme).
\definecolor{earthDark}{rgb}{0.30, 0.40, 0.18}
\definecolor{earthSienna}{rgb}{0.55, 0.34, 0.20}

\def\wacvPaperID{1157}
\def\confName{WACV}
\def\confYear{2027}

% Float settings mirror main.tex so appendix tables lay out identically.
\renewcommand{\topfraction}{0.95}
\renewcommand{\dbltopfraction}{0.95}
\renewcommand{\textfraction}{0.05}
\renewcommand{\floatpagefraction}{0.5}
\renewcommand{\dblfloatpagefraction}{0.5}
\setcounter{topnumber}{2}
\setcounter{dbltopnumber}{2}

\begin{document}

% \title defines \thetitle, which \maketitlesupplementary prints. No \maketitle.
\title{Fields of the Planet: Field Boundary Mapping Beyond 10m}

% One-column for the full-width appendix tables. Title banner is set manually
% (rather than \maketitlesupplementary) so it flows into the content instead of
% leaving a near-empty two-column title page.
\onecolumn
\begin{center}
\Large
\textbf{\thetitle}\\
\vspace{0.5em}Supplementary Material \\
\vspace{1.0em}
\end{center}
\appendix
% Continue table numbering from the main paper (main ends at Table 3).
\setcounter{table}{3}

\section{Polygon-metric definitions}
\label{app:metrics}

All object-level metrics are computed on vector geometries for the field-interior class. Predicted masks are vectorized into polygons and scored against the original FTW ground-truth polygons; pixel IoU is computed before vectorization and reported only for continuity with the released PRUE protocol.

\paragraph{Vectorization and matching.}
For each patch, the post-processed field-interior prediction is converted into polygons $\mathcal{P}=\{p_1,\ldots,p_M\}$ and compared with the ground-truth polygons $\mathcal{G}=\{g_1,\ldots,g_N\}$. All intersections and unions use polygon areas from the vector geometries. At IoU threshold $\tau$, a predicted polygon and a ground-truth polygon are matched when
\begin{equation}
\mathrm{IoU}(g,p)=
\frac{\mathrm{area}(g\cap p)}
{\mathrm{area}(g\cup p)}
>\tau .
\end{equation}
For $\tau\ge0.5$, disjoint polygon interiors make the match set $\mathcal{M}_\tau$ a partial one-to-one matching. We then define $\mathrm{TP}_\tau=|\mathcal{M}_\tau|$, $\mathrm{FP}_\tau=M-\mathrm{TP}_\tau$, and $\mathrm{FN}_\tau=N-\mathrm{TP}_\tau$.

\paragraph{Object F1 and panoptic quality.}
Object F1 is reported at IoU$=0.5$ and over the COCO threshold grid $\mathcal{T}={0.5,0.55,\ldots,0.95}$~\cite{lin2014coco}:
\begin{equation}
\mathrm{F1}_{[.5:.95]}=
\frac{1}{|\mathcal{T}|}
\sum_{\tau\in\mathcal{T}}\mathrm{F1}_{\tau}.
\end{equation}
Panoptic quality follows Kirillov et al.~\cite{kirillov2019panoptic}. With matches fixed at $\tau=0.5$, recognition quality is $\mathrm{RQ}=\mathrm{F1}_{0.5}$, segmentation quality is the mean IoU of matched polygons,
\begin{equation}
\mathrm{SQ}=
\frac{1}{\mathrm{TP}_{0.5}}
\sum_{(g,p)\in\mathcal{M}_{0.5}}
\mathrm{IoU}(g,p),
\end{equation}
and $\mathrm{PQ}=\mathrm{SQ}\times\mathrm{RQ}$. If $\mathrm{TP}_{0.5}=0$, we set $\mathrm{SQ}=0$ and $\mathrm{PQ}=0$; matched-boundary error is undefined.

\paragraph{Polygon-count error.}
Let $\bar{M}$ and $\bar{N}$ denote the mean numbers of predicted and ground-truth polygons per patch within a country. We report the normalized count discrepancy
\begin{equation}
|\Delta N|/N=
\frac{|\bar{M}-\bar{N}|}{\bar{N}},
\end{equation}
computed per country and then macro-averaged. This is a coarse object-count check independent of IoU matching; it measures count mismatch magnitude but not whether errors come from over- or under-prediction.

\paragraph{Boundary error.}
For each matched pair $(g,p)\in\mathcal{M}_{0.5}$, we compute a symmetric boundary chamfer distance between the one-pixel-wide ground-truth and predicted boundary sets on the evaluated raster grid. Nearest-boundary distances are read from Euclidean distance transforms, averaged in both directions, and converted to meters by multiplying by the evaluated grid's ground sample distance: 3m for PlanetScope and 10m for native Sentinel-2. We report the mean and 95th percentile over matched pairs. Boundary error therefore measures localization for fields that were recovered; missed and hallucinated fields are penalized through RQ and object F1.

\paragraph{Pixel IoU.}
For continuity with PRUE~\cite{muhawenayo2026prue}, we also report the pixel-level Jaccard index for the field-interior class,
\begin{equation}
\mathrm{IoU}_{\mathrm{px}} =
\frac{|\hat{Y}_{\mathrm{field}}\cap Y_{\mathrm{field}}|}
{|\hat{Y}_{\mathrm{field}}\cup Y_{\mathrm{field}}|},
\end{equation}
computed on raster masks before vectorization. Unlike the polygon metrics, pixel IoU has no notion of individual parcels.

\paragraph{Aggregation.}
Metrics are first pooled within each country: $\mathrm{TP}/\mathrm{FP}/\mathrm{FN}$ counts and matched IoUs are accumulated across patches before computing RQ, SQ, PQ, and $\mathrm{F1}_{[.5:.95]}$, while chamfer is averaged over that country's matched pairs. Main-paper Table~1 reports macro-averages over the ten dense-label held-out countries. If a configuration produces no matched polygons in a country, boundary error is undefined for that country and excluded from the boundary-error average; all other columns remain defined.

\paragraph{Field-area breakdown.}
For main-paper Table~3, ground-truth polygons are assigned by UTM area $A$ to \emph{small} ($A<0.5$\,ha), \emph{medium} ($0.5\le A<2$\,ha), or \emph{large} ($A\ge2$\,ha). The 2\,ha threshold follows the standard smallholder ceiling~\cite{lowder2016farms}, while 0.5\,ha separates the dominant sub-hectare fields in FTW. Matched true positives and false negatives are binned by ground-truth area; unmatched predictions are counted as false positives in the bin determined by their predicted area. PQ is reported per bin and macro-averaged over the ten dense-label held-out countries.


\section{Upsampled Sentinel-2 Control}
\label{app:upsampled_s2}

This control separates output-grid resolution from image resolution. We give the Sentinel-2 model a finer 512-pixel prediction grid without adding new image information, either by upsampling at evaluation time or by training on bilinearly upsampled Sentinel-2 imagery. Both settings improve object-level polygon scores, but neither recovers the meter-scale boundary precision of native 3m PlanetScope imagery.

\begin{table}[h]
    \centering
    \small
    \setlength{\tabcolsep}{5pt}
    \caption{\textbf{Upsampling Sentinel-2 improves object recovery but not boundary precision.} Results use B3 models on the dense held-out split with WS + TTA. Rows share the same upsampled-comparable scoring protocol.}
    \label{tab:upsampled_s2_main}
    \begin{tabular}{lcccc}
    \toprule
    Sentinel-2 grid & PQ & SQ & RQ & Bd.\ err (m) \\
    \midrule
    Native ($256$, $10$m)         & $23.8$ & $67.8$ & $30.8$ & $18.6$ \\
    Upsampled at eval ($512$)     & $27.4$ & $75.8$ & $34.9$ & $31.3$ \\
    Trained upsampled ($512$)     & $31.5$ & $77.5$ & $39.5$ & $28.0$ \\
    \midrule
    Planet ($512$, native $3$m)   & $\mathbf{36.0}$ & $\mathbf{77.9}$ & $\mathbf{45.2}$ & $\mathbf{7.4}$ \\
    \bottomrule
    \end{tabular}
\end{table}

\section{Dataset scope and per-tile UDM2 quality}
\label{app:scope}
\label{app:udm2_country}

Per-region usable-patch rates (clear$\,\ge\,$95\%, unusable$\,\le\,$5\%) vary widely across the release; smallholder regions (e.g.\ India) have low \emph{coverage} per patch ($\sim$1\% field pixels) but high clear-sky fractions. The full per-region index ships with the release as a CSV.

\begin{table}[h]
\centering
\begin{minipage}[t]{0.54\linewidth}
\centering
\captionof{table}{\textbf{Dataset scope.}
FTP pairs FTW's 24 countries / 25 labeled regions with two seasonal PlanetScope windows per patch. \emph{Success rate} denotes the fraction of FTW patch-window targets with a successfully extracted PlanetScope image.}
\label{tab:scope}
\small
\begin{tabular}{lr}
\toprule
\textbf{Property} & \textbf{Value} \\
\midrule
\multicolumn{2}{l}{\cellcolor{black!7}\textbf{Geographic coverage}}\\
Countries / regions / continents & 24\,/\,25\,/\,4 \\
FTW patches                  & 70{,}484 \\
(Patch, window) targets      & 140{,}968 \\
\midrule
\multicolumn{2}{l}{\cellcolor{black!7}\textbf{Source imagery}}\\
Unique PSScene COGs touched  & 6{,}113 \\
Mean windows per scene       & 21.8 \\
\midrule
\multicolumn{2}{l}{\cellcolor{black!7}\textbf{Released artifacts}}\\
PlanetScope SR (4-band, uint16) & 133{,}168 \\
Per-window UDM2 QA statistics   & 129{,}490 \\
3-class labels (uint8, 1/patch) & 66{,}584 \\
\midrule
\multicolumn{2}{l}{\cellcolor{black!7}\textbf{Yield}}\\
Successfully matched pairs      & 133{,}168 \\
\textbf{Success rate}           & \textbf{94.5\%} \\
Total dataset size              & 102\,GB \\
\bottomrule
\end{tabular}
\end{minipage}\hfill
\begin{minipage}[t]{0.43\linewidth}
\centering
\captionof{table}{\textbf{Per-tile UDM2 coverage.}
Percentage of patch-window tiles whose UDM2 class coverage exceeds each threshold. Most tiles are clear ($95.9\%$ are $\ge 50\%$ clear); tiles with $\ge 50\%$ cloud, shadow, haze, or unusable pixels are resampling candidates.}
\label{tab:udm2}
\small
\begin{tabular}{l S[table-format=2.1] S[table-format=2.1] S[table-format=2.1]}
\toprule
{Band} & {any (\textgreater0)} & {\bfseries $\ge 50\%$} & {$\ge 90\%$} \\
\midrule
\texttt{clear}      & 98.9 & \bfseries 95.9 & 91.1 \\
\texttt{cloud}      &  9.0 & \bfseries  1.0 &  0.2 \\
\texttt{shadow}     &  6.9 & \bfseries  0.2 &  0.0 \\
\texttt{light\_haze}&  5.1 & \bfseries  2.0 &  1.1 \\
\texttt{unusable}   & 11.3 & \bfseries  1.3 &  0.4 \\
\bottomrule
\end{tabular}
\par\vspace{2pt}\footnotesize \% of (patch, window) tiles; each tile's class area as a fraction of the window.
\end{minipage}
\end{table}

\section{Held-out post-processing sweep}
\label{app:heldout_sweep}

\Cref{tab:heldout} reports the full watershed/D4-TTA sweep for PRUE+ models on the ten-country dense held-out set. The released FTW-PRUE rows are included as external reference values from the original PRUE test split, not as held-out macro-averages.

\paragraph{PRUE+ augmentation settings.}
PRUE+ is implemented with Kornia~\cite{riba2020kornia}. It keeps the released PRUE preprocessing augmentations---\texttt{preprocess\_aug} with divisor sampled from $[5000,15000]$ and \texttt{RandomResizedCrop} with scale $0.3$--$0.9$, ratio $0.75$--$1.33$, and $p=0.5$---then adds seasonal-window swap and per-band gamma jitter ($\gamma\sim\mathcal{U}[0.8,1.2]$, $p=0.3$). The final geometry/noise bundle adds affine intensity jitter ($a\sim\mathcal{U}[0.9,1.1]$, $b\sim\mathcal{U}[-0.02,0.02]$, $p=0.3$), Gaussian blur ($3\times3$, $\sigma\sim\mathcal{U}[0.5,1.5]$, $p=0.3$), Gaussian noise ($\sigma=0.015$, $p=0.3$), single-window dropout ($p=0.15$), rotation ($\pm30^\circ$, $p=0.5$), shear ($\pm5^\circ$, $p=0.3$), and boundary-thickness jitter by random binary dilation of $0$--$2$ pixels ($p=0.5$).

\begin{table*}[h]
    \centering
    \scriptsize
    \caption{\textbf{Held-out post-processing sweep.}
    Object F1 for B3/B7 PRUE+ models across watershed (WS) and D4 test-time augmentation (TTA) settings. Dense-label macro-averages exclude presence-only Kenya and retain Portugal; released FTW-PRUE rows are published PRUE reference values. Best per row is bolded.}
    \label{tab:heldout}
    \begin{tabular}{lcccc}
\toprule
Model & Pix IoU & Pix Prec & Pix Rec & Obj F1 \\
\midrule
\multicolumn{5}{l}{\textit{9 (excl K+P) held-out countries (macro avg)}}\\
S2 CCBY (B3, CE) \cite{kerner2024ftw} & \textbf{0.608} & \textbf{0.834} & 0.706 & 0.317 \\
FTW-Planet B3 (PRUE recipe) & 0.541 & 0.815 & 0.637 & 0.323 \\
FTW-Planet B3 + SDF aux & 0.520 & 0.814 & 0.616 & 0.320 \\
FTW-Planet B3 + frame-field aux & 0.542 & 0.825 & 0.632 & 0.302 \\
FTW-Planet B7 + SDF aux & 0.571 & 0.785 & \textbf{0.709} & \textbf{0.334} \\
\midrule
\multicolumn{5}{l}{\textit{all 11 held-out countries (macro avg)}}\\
S2 CCBY (B3, CE) \cite{kerner2024ftw} & \textbf{0.537} & \textbf{0.828} & 0.619 & 0.262 \\
FTW-Planet B3 (PRUE recipe) & 0.444 & 0.750 & 0.570 & 0.265 \\
FTW-Planet B3 + SDF aux & 0.478 & 0.734 & 0.623 & \textbf{0.279} \\
FTW-Planet B3 + frame-field aux & 0.446 & 0.700 & 0.574 & 0.249 \\
FTW-Planet B7 + SDF aux & 0.468 & 0.642 & \textbf{0.641} & 0.274 \\
\bottomrule
\end{tabular}

\end{table*}

\begin{table}[h]
    \centering
    \footnotesize
    \setlength{\tabcolsep}{3pt}
    \caption{\textbf{Recipe ablation summary.}
    Object-F1 changes on the ten-country dense held-out set, using WS + TTA unless noted. Checkmarks indicate components kept in FTP-PRUE+; $\times$ marks rejected variants.}
    \label{tab:ablation_summary}
    \begin{tabular}{@{}p{0.62\linewidth} @{\hspace{2pt}}c@{\hspace{2pt}} r@{}}
    \toprule
    \textbf{Lever} & \textbf{Use} & \textbf{$\Delta$ ObjF1} \\
    \midrule
    \multicolumn{3}{@{}l}{\cellcolor{black!7}\textbf{Augmentation \& training}} \\
    PRUE preprocess+resize aug.~\cite{muhawenayo2026prue} & \checkmark & {\color{earthDark}\bfseries +2.8} \\
    \texttt{swap\_order} + per-band $\gamma$ & \checkmark & {\color{earthDark}\bfseries +1.0} \\
    Geometry/noise bundle (PRUE+: affine, blur, noise, window dropout, rot, shear, jitter) & \checkmark & {\color{earthDark}\bfseries +4.5} \\
    Backbone B3 $\to$ B7 (no gain, larger) & $\times$ & {\color{black!50}\bfseries $\sim$0} \\
    Soft clDice~\cite{shit2021cldice} on boundary & $\times$ & {\color{earthSienna}\bfseries failed} \\
    SDF auxiliary head~\cite{hayder2017bais,bischke2019buildings} (with PRUE+) & $\times$ & {\color{earthSienna}\bfseries -3.1} \\
    Frame-field head~\cite{girard2021ffl} & $\times$ & {\color{earthSienna}\bfseries -1.7} \\
    CutMix~\cite{yun2019cutmix} (2\,px ignore) & $\times$ & {\color{earthSienna}\bfseries -0.1} \\
    \midrule
    \multicolumn{3}{@{}l}{\cellcolor{black!7}\textbf{Inference \& post-processing}} \\
    Marker-controlled watershed & \checkmark & {\color{earthDark}\bfseries +0.5} \\
    D4 8-way TTA (with PRUE+) & \checkmark & {\color{earthDark}\bfseries +0.9} \\
    \midrule
    \multicolumn{3}{@{}l}{\cellcolor{black!7}\textbf{Checkpoint selection}} \\
    Last vs best-by-val checkpoint & either & {\color{black!50}\bfseries $\le0.3$} \\
    \bottomrule
    \end{tabular}
\end{table}

{\small
\bibliographystyle{ieeenat_fullname}
\bibliography{refs}
}

\end{document}
